\documentclass[11pt]{article}
\usepackage[T1]{fontenc}
\usepackage[a4paper,margin=27mm]{geometry}
\usepackage{lmodern,amsmath,amssymb,amsthm,booktabs,microtype}
\usepackage[round,authoryear]{natbib}
\usepackage{xurl,hyperref}
\hypersetup{hidelinks}
\newtheorem{theorem}{Theorem}[section]
\newtheorem{lemma}[theorem]{Lemma}
\newtheorem{proposition}[theorem]{Proposition}
\newtheorem{corollary}[theorem]{Corollary}
\theoremstyle{remark}\newtheorem{remark}[theorem]{Remark}
\newcommand{\Q}{\mathbb Q}
\newcommand{\Z}{\mathbb Z}
\newcommand{\F}{\mathbb F}
\newcommand{\Gal}{\operatorname{Gal}}
\newcommand{\rd}{\operatorname{rd}}
\newcommand{\cond}{\operatorname{cond}}
\newcommand{\N}{\operatorname N}
\newcommand{\tr}{\operatorname{tr}}
\newcommand{\workbenchurl}{https://github.com/KokunoYumeto/zeta-and-l-functions-workbench}
\newcommand{\recordurl}{https://github.com/KokunoYumeto/mathematics-workbenches}
\title{CM hosts and prime transfer in Alp\"oge's theta-lift hypotheses}
\author{}
\date{24 November 2026}
\begin{document}
\maketitle

\begin{abstract}
Alp\"oge's theta-lift formalism imposes arithmetic conditions on its CM
host in addition to its numerical size conditions. We prove that every
literal host contains a maximal real cyclotomic field whose
conductor exceeds twice the comparison scale. Consequently no odd
prime ideal has norm in that range, and all corresponding incomplete
Hecke logarithmic-derivative coefficients vanish. This unconditional
obstruction also applies to arbitrary signed mixtures of such ordinary
host statistics. It does not apply to the formalism's different local
gain, which is indexed by rational residue characteristic.
For the cubic character of $\Q(\omega,\sqrt[3]{269})/\Q(\omega)$,
we construct relaxed growing CM fields that preserve every Euler
factor in the range and satisfy all five isolated numerical conditions.
Under Dirichlet GRH we give a uniform discriminant lower bound and a
lower bound for the coefficient mass of relaxed weighted transfers.
A fixed quotient containing the cubic target factor has an exact
Artin projector with an explicit integral cost. The source's separated
highest weights prevent the naive deck action on a single coefficient
module. Uniform lattice gains pass to compatible projections, but a
coherent equivariant theta realization remains to be constructed.
No new Riemann hypothesis is proved.
The text and finite computations were produced by GPT-6 Astra in
ChatGPT Pro mode; the recorded checks are by the same model, and no
human has checked the text.
\end{abstract}

\section{Context and main statements}

We examine a precise question about the host field in the formalism
accompanying \citet{Alpoge2026}: can growth of the degree be reconciled
with an exact comparison of the cubic prime statistic?
Class-field theory supplies such a comparison for relaxed hosts.
The literal source conditions instead force an empty range of small
odd prime-ideal norms, by an unconditional algebraic argument.
For the broader abelian class we also state an analytic obstruction
under Dirichlet GRH, isolating the degree-one input of Alp\"oge's
theorem; no nonabelian Riemann hypothesis enters.

Put
\[
 K=\Q(\omega),\qquad E=K(\alpha),\qquad \alpha^3=269,
 \qquad \omega^2+\omega+1=0.
\]
Let $\chi$ be either faithful character of $\Gal(E/K)$ and let
$\rho=\operatorname{Ind}_{G_K}^{G_\Q}\chi$ be the two-dimensional
irreducible representation of $\Gal(E/\Q)\simeq S_3$.
For $p\nmid3\cdot269$, write
\[
 a_{p,k}=\tr\rho(\operatorname{Frob}_p^k),\qquad a_p=a_{p,1}.
\]
At a prime $p\equiv1\pmod3$, $a_p$ is $2$ or $-1$; at a prime
$p\equiv2\pmod3$, it is zero.
For an abelian field $F\supset K$, put $h=[F:K]$ and
$\chi_F=\chi\circ\N_{F/K}$, interpreted as the base-changed
finite-order Hecke character. Let $L^S$ omit the rational primes
$S=\{3,269\}$, and define coefficients by
\[
 -\frac{(L_F^S)'(s,\chi_F)}{L_F^S(s,\chi_F)}
       =\sum_{n\ge1}\frac{b_F(n)}{n^s}
       \qquad(\Re s>1).
                                                        \tag{1.1}
\]
For $K$, the corresponding coefficient at $p^k$ is
$b_K(p^k)=a_{p,k}\log p$ if $p\notin S$ and is zero otherwise.

\begin{theorem}[No small odd prime-ideal norms in literal hosts]
\label{thm:empty-norms}
Let $F$ be a literal host of \citet{AlpogeLean2026}, with its frame
parameter $e$, degree $[F:\Q]=4n$, and $\texttt{Host.o\_pow}$
condition. Then
\[
 B_e=\Q(\zeta_{2^{e+2}})^+\subset F,\qquad
 n\le2^{65e/64}.
\]
If $x>1$, $0<\epsilon<1/52$, and its lower degree condition
$n\ge x^{2+8\epsilon-\epsilon/100}$ holds, then
$2^{e+2}>2x+1$. In every finite extension $F'/F$, no prime
ideal of odd residue characteristic has norm at most $2x$.
For every finite-order Hecke character of $F'$, all coefficients
through $2x$ of its logarithmic derivative, after deleting the
Euler factors above $2$, are therefore zero. Every finite complex
linear combination of these ordinary coefficient sequences is
zero in that range. In particular, for $x\ge7/2$ it cannot equal
the cubic target sequence with the factors at $2,3,269$ deleted.
All these assertions are unconditional. They concern ordinary
Hecke coefficients ordered by absolute norm, not the source's
local theta gains ordered by rational residue characteristic.
For each odd unramified prime $p\le2x$, the order of the
$A$-coordinate of its Frobenius is at least
$2^{e+2}/(p+1)>x^{\epsilon/4}$; thus the source's long-orbit
condition holds on this range even though its small-norm
Hecke coefficients vanish.
For fixed $\epsilon$ and fixed local constant $\texttt{kII}$, if the upper
\texttt{SizeK} bound also holds, then every such prime in the
source's \texttt{inertClass} satisfies \texttt{Cand},
$\Sigma\le x^{-\epsilon}$, and the negation of \texttt{NoGain},
for all sufficiently large scales. Thus its served-orbit set equals
its good-orbit set for every word. This does not assert that the
good-orbit set is nonempty or identify a cubic gain.
\end{theorem}

\begin{theorem}[Common hosts with exact transfer]\label{thm:construction}
Fix $0<\epsilon<1/52$. For every sufficiently large real $x$ there
is a CM field $F_x$, abelian over $\Q$ and containing $K$, such that:
\begin{enumerate}
\item $[F_x:\Q]=4n$ for a power of two $n$, with
\[
 \tfrac12x^{2+8\epsilon}<n\le x^{2+8\epsilon}.
                                                        \tag{1.2}
\]
In particular $[F_x:\Q]>4x^2$ for sufficiently large $x$.
\item Every rational prime $p\le2x$ split in $K$ splits completely
in $F_x$. Every prime $p\equiv2\pmod3$ with $p\le2x$ has residue
degree exactly two in $F_x$. Every prime of $K$ above such a rational
prime splits completely in $F_x/K$.
\item The character $\chi_{F_x}$ is nontrivial,
\[
 \mathfrak f(\chi_{F_x})=269\mathcal O_{F_x},\qquad
 \N\mathfrak f(\chi_{F_x})=269^{[F_x:\Q]},
                                                        \tag{1.3}
\]
and, for every $p\le2x$, $p\notin S$, its local Euler factor is
\[
 L_{F_x,p}(s,\chi_{F_x})=L_{K,p}(s,\chi)^h.
                                                        \tag{1.4}
\]
Consequently $b_{F_x}(n)/h=b_K(n)$ for every integer $n\le2x$.
\item There are distinct primes $\ell_j\equiv1\pmod4$ in
$(2x,10x)$, and a subset $J$ of their indices, such that
\[
 \rd(F_x)=\sqrt{3\prod_{j\in J}\ell_j},\qquad
 \log\rd(F_x)=O(x).
                                                        \tag{1.5}
\]
The implied constant is absolute after $\epsilon$ is restricted as
above and $x$ exceeds a suitable threshold.
\item The degree in (1.2) can be accompanied by integers $m,k$ and
positive numbers $g_0,g_1$ satisfying every inequality (3.1) below,
including both upper bounds on $n$ and $k$ and all the displayed
rank conditions. The $g_i$ can be chosen to be distinct primes.
\end{enumerate}
\end{theorem}

The next theorem applies more broadly when the literal cyclotomic
condition is relaxed. A degree-one prime means degree one over $\Q$.
At primes where $\chi_F$ is ramified its local character value is
understood to be zero.

\begin{theorem}[Discriminant cost of exact prime transfer]
\label{thm:discriminant}
Assume Dirichlet GRH. There are absolute constants $c>0$ and $x_0$
such that the following holds. Let $x\ge x_0$, and let $F/\Q$
be a finite abelian extension with $K\subsetneq F$. If, for every
prime $p\le2x$ with $p\equiv1\pmod3$,
\[
 \frac1{[F:K]}
   \sum_{\substack{\mathfrak P\mid p\text{ in }F\\
                    f(\mathfrak P/p)=1}}
                \chi_F(\mathfrak P)=a_p,
                                                        \tag{1.6}
\]
then every such $p$ splits completely in $F$, and
\[
                         \log\rd(F)\ge c\sqrt{x}.
                                                        \tag{1.7}
\]
The constants do not depend on $F$, its degree, or its conductor.
\end{theorem}

\begin{theorem}[Incompatibility with the literal host conditions]
\label{thm:literal}
Assume Dirichlet GRH. Fix $0<\epsilon<1/52$. For all sufficiently
large $x$, no host in the literal formalism of \citet{AlpogeLean2026}
can both contain $K$ and satisfy (1.6), together with its
\texttt{SizeT}, \texttt{SizeCond}, and \texttt{SizeQD} conditions.
More generally, the same conclusion holds for any proper abelian
extension $F/K$ over $\Q$ satisfying (1.6) and the scalar conditions
\[
 Q\ge1,\quad q,d\ge1,\quad t\ge1,\qquad
 (Qqd)^{36}t^{150}\le e^{x^\epsilon/2},\qquad
 \rd(F)\le\frac Q2\sqrt{qd}.
                                                        \tag{1.8}
\]
Independently of GRH and of $x$, the fields constructed in
Theorem~\ref{thm:construction} cannot be literal hosts: their Galois
groups have exponent two, whereas a literal host has a cyclic
quotient of order four.
\end{theorem}

\begin{theorem}[Cost of relaxed signed mixtures]\label{thm:mixture-cost}
Assume Dirichlet GRH. Fix $0<\epsilon<1/52$, $c>0$, and a nonnegative
nonzero $W\in C_c^\infty((1,2))$. At each scale let $F_i/\Q$ be
finitely many abelian fields containing $K$, with
$[F_i:\Q]\ge c x^2$ and $\log\rd(F_i)\le x^\epsilon/72$.
At a prime $p\equiv1\pmod3$ define $\alpha_i(p)=0$ if the
absolute residue degree exceeds one, and $\alpha_i(p)=1/e_i(p)$
otherwise. If complex numbers $c_i$ satisfy
$\sum_i c_i\alpha_i(p)=1$ for every such prime with $W(p/x)>0$,
then, for sufficiently large $x$,
\[
                   \sum_i|c_i|\gg_{W,c,\epsilon}x^{1/2-\epsilon}.
\]
The same conclusion holds when the weighted absolute defect
\[
 \sum_{p\equiv1(3)}\left|1-\sum_i c_i\alpha_i(p)\right|
                       \log p\,W(p/x)
\]
is at most one quarter of
$\sum_{p\equiv1(3)}\log p\,W(p/x)$.
The bound is uniform in the fields and the number of summands.
\end{theorem}

\begin{theorem}[A fixed quotient and its projector]\label{thm:projector}
Let $B=\Q(\zeta_{16})^+$, $F_0=KB$, and let $F/\Q$ be any
finite abelian extension containing $F_0$. Set $G=\Gal(F/K)$,
$H=\Gal(F/F_0)$, and $r=|H|$. On
$V=\operatorname{Ind}_{G_F}^{G_K}(\chi|_{G_F})
\simeq\chi\otimes\mathbb C[G]$, the average of the right $H$-actions
is an orthogonal projector of norm one and rank four. Its image has
$L$-function
\[
 J_0(s)=L_{F_0}(s,\chi_{F_0})
       =\prod_{\psi\in\widehat{\Gal(F_0/K)}}L_K(s,\chi\psi).
\]
Its projected Frobenius traces give the coefficients of this fixed
function at every prime power, using inertia invariants at ramified
primes. All four factors are entire, and RH for $J_0$ implies RH
for $L_K(s,\chi)$. For $\mathcal O=\Z[\omega]$, the quotient of
the projected regular lattice by its invariant lattice is exactly
$(\mathcal O/r\mathcal O)^4$, of cardinality $r^8$.
This is an Artin-representation statement, not a theta realization.
\end{theorem}

\begin{theorem}[Noninvariance of one coefficient module]
\label{thm:highest}
Fix one literal family, a fixed $\epsilon>0$, and its
scale-independent cloud constant $C_c$. Assume its members satisfy
\texttt{SizeTT} and \texttt{SizeDel}.
For all sufficiently large scales, let $G=\Gal(F/\Q)$, let
$\Psi\subset G$ be its CM-label subgroup, and let $H\subset G$
have index at most eight. Then the layer-zero highest-weight tuple
of every word member is not invariant under multiplication by
$H\cap\Psi$ on $\Psi$. Its external tensor-product coefficient
representation admits no action of all of $H\cap\Psi$ that
implements coordinate permutation and intertwines the corresponding
permutation of the product $\mathrm{GL}_k$.
This assertion concerns one fixed labelled representation. It does
not exclude a direct sum of conjugated representations, a relabelling
of different transported objects, or another action with additional
structure.
\end{theorem}

\section*{Proof overview}
The source's Artin map identifies a specific cyclotomic quotient.
Galois correspondence turns it into a real cyclotomic subfield.
A second source condition bounds the index of a distinguished cyclic
subgroup; together these force the cyclotomic conductor beyond the
entire prime-ideal norm range. This proves the strongest literal
obstruction without an analytic hypothesis. The distinction between
absolute norm and rational residue characteristic explains why it
does not contradict the source's local theta mechanism.

For relaxed hosts we start with more real quadratic generators than
desired and quotient by the Frobenius elements of all split primes
up to $2x$. Finite linear algebra leaves the desired degree. Since
this Galois group has exponent two, inert rational primes also split
over their degree-two primes in $K$. This gives exact Euler factors.
Conversely, exact pointwise transfer in an abelian field forces
complete splitting. A uniform smooth Dirichlet explicit formula and
the conductor--discriminant formula give the discriminant cost.
The same normalized explicit formula, applied to the nonnegative
Dedekind-zeta coefficients, bounds how much prime mass one relaxed
host can supply and hence bounds the coefficient mass of a mixture.

A different comparison keeps one fixed quotient and projects the
regular Artin representation onto it. Its four entire factors include
the original cubic function. The projector and its lattice index are
explicit. To test whether this identity can enter the theta
chain, we use the source's highest weights: the small top set has
weights exponentially larger than the remaining ones and cannot be
invariant under the proposed large deck group. This prevents the
naive internal action on one coefficient module. Uniform lattice
inclusions do pass through compatible equivariant projections; the
remaining construction must supply those actions and lattices and the
correct projected trace and centering identity.

\section{The real cyclotomic subfield and the empty norm range}
\label{sec:literal-small}

We first prove the unconditional conclusion, stating the precise source
data that enter it. Write $F$ for the field called the host field in
\citet{AlpogeLean2026}. A frame gives a finite abelian group $A$ of
order $n$, a surjection
\[
 h:A\longrightarrow H_e=(\Z/2^{e+2}\Z)^\times/\{\pm1\},
 \qquad a_0\in A,
                                                        \tag{2.1}
\]
and the implication $h(ja_0)=0\Rightarrow ja_0=0$ for every integer
$j$. Put $n_0=\operatorname{ord}(a_0)$ and $o=n/n_0$.
The full host has $e\ge2$, $o^{64}\le2^e$, and
$\Gal(F/\Q)\simeq C_2^2\times A$.
These are the source declarations \texttt{Frame}, \texttt{Ranks.n},
and \texttt{Host.o\_pow}; the order $n$ is the same parameter as
in its degree window.

There is more information than an abstract cyclic quotient. The
host includes an embedding
\[
 \iota:F\hookrightarrow C=\Q(\zeta_N),\qquad
 N=Q_Sqd,\qquad 2^{e+2}\mid Q_S,
                                                        \tag{2.2}
\]
and an Artin map from $(\Z/N\Z)^\times$ onto $\Gal(F/\Q)$.
The source equality \texttt{art\_spec} says that this map is the
restriction of the cyclotomic action through $\iota$.
The equality \texttt{art\_h} identifies its projection to $H_e$
with reduction modulo $2^{e+2}$ followed by quotienting by $\{\pm1\}$.
We identify $F$ with its image under $\iota$.

\begin{lemma}[A common subfield]\label{lem:actual-cyclotomic}
Every such host contains
\[
       B_e=\Q(\zeta_{2^{e+2}})^+,
       \qquad B=\Q(\zeta_{16})^+\subset B_e.
                                                        \tag{2.3}
\]
\end{lemma}
\begin{proof}
An automorphism of $C$ is given by a unit $u$ modulo $N$.
If it fixes $F$, the restriction equality in (2.2) and injectivity
of $\iota$ make its Artin image the identity. Its projection to
$H_e$ is therefore the identity. By \texttt{art\_h}, this means
$u\equiv\pm1\pmod{2^{e+2}}$, exactly the condition that the
automorphism fix $B_e$. Hence
$\Gal(C/F)\subset\Gal(C/B_e)$, and the Galois correspondence
gives $B_e\subset F$. The second inclusion follows from $e\ge2$.
\end{proof}

\begin{lemma}[Degree forces a large cyclotomic conductor]
\label{lem:large-tower}
For the full host,
\[
 n\le2^{65e/64},\qquad
 M:=2^{e+2}\ge4n^{64/65}.
                                                        \tag{2.4}
\]
If $x>1$, $\epsilon>0$, and
$n\ge x^{2+8\epsilon-\epsilon/100}$, then $M>2x+1$.
\end{lemma}
\begin{proof}
The implication following (2.1) makes $h$ injective on
$\langle a_0\rangle$: apply it to the difference of any two
elements with equal image. The group $H_e$ has order $2^e$, so
$n_0\le2^e$. Also $n_0\mid n$ by Lagrange's theorem and thus
$n=o n_0$. Combining this with $o^{64}\le2^e$ proves (2.4).
For completeness, $|H_e|=2^e$ because the group of units modulo
$2^{e+2}$ has order $2^{e+1}$ and $-1$ has order two.
The lower degree window now gives
\[
 M\ge4x^{(64/65)(2+799\epsilon/100)}>4x>2x+1.
                                                        \tag{2.5}
\]
The exponent exceeds one, and the final inequality uses $x>1$.
\end{proof}

\begin{proof}[Proof of Theorem~\ref{thm:empty-norms}]
Let $\mathfrak P$ be a prime of any extension $F'/F$ above an odd
rational prime $p$. Write $\N\mathfrak P=p^f$. Its restriction
to $B_e$ has residue degree dividing $f$, since the corresponding
residue field embeds in that at $\mathfrak P$. The field $B_e$
is unramified at $p$; its Frobenius is the class of $p$ modulo
$M$, up to sign. Consequently
\[
                         p^f\equiv\pm1\pmod M.
                                                        \tag{2.6}
\]
If $p^f\le2x$, (2.5) gives $1<p^f<M-1$, contradicting (2.6).
This argument does not assume that $F$ or $F'$ is unramified at $p$.

The logarithmic derivative of a Hecke Euler product is a sum over
powers of prime-ideal norms. A coefficient of index at most $2x$
and odd residue characteristic would therefore require a prime
ideal of norm at most $2x$, which we have excluded. Deleting the
factors over $2$ proves the coefficientwise vanishing assertion.
All its finite linear combinations vanish as well. At $p=7$ the
cubic trace is $-1$, because $269\equiv3\pmod7$ is not a cube.
Thus once $x\ge7/2$ these coefficients cannot recover all the
target coefficients through $2x$.
\end{proof}

The exclusion is stronger than the fixed common-subfield obstruction.
Even without the degree window, (2.3) implies that every odd prime
$p\not\equiv\pm1\pmod{16}$ has no degree-one prime in a literal
host. Among split-$K$ classes modulo $48$, the only potentially
visible classes are $1,31$; the missing classes are
$7,13,19,25,37,43$. With the full degree window every odd prime
in the range is missing, including the two formerly visible classes.

\subsection{Why this is not an obstruction to the source's theta gain}

There is also a positive consequence for the local orbit mechanism.
For an odd prime $p\le2x$, let
\[
 v=\begin{cases}v_2(p-1),&p\equiv1\pmod4,\\
                 v_2(p+1),&p\equiv3\pmod4.
   \end{cases}
\]
Choose $u=p$ or $u=-p$ so that $u\equiv1\pmod4$. Then
$v_2(u-1)=v\ge2$, and successive factorizations of
$u^{2^{j+1}}-1$ show
$v_2(u^{2^j}-1)=v+j$: the second factor always has valuation one.
A power of $u$ is $1$ modulo four, so it cannot be $-1$ modulo
$M\ge16$. As $M>p+1$, the class of $p$ in $H_e$ therefore has
order exactly
\[
 \frac M{2^v}\ge\frac M{p+1}\ge\frac M{2x+1}
   >\frac43x^{\beta-1}>x^{\epsilon/4},\qquad
 \beta=\frac{64}{65}\left(2+\frac{799\epsilon}{100}\right).
                                                        \tag{2.7}
\]
Here $\beta-1>\epsilon/4$. At an unramified prime the source
Artin equality identifies this class with the image of its
$A$-coordinate, whose order is consequently at least as large.
This proves the long-orbit assertion of
Theorem~\ref{thm:empty-norms}. Its order is even, so the
$A$-coordinate has positive even order. The source predicate
\texttt{HEv} holds, and both its zero-coordinate and odd-order
\texttt{IsCs} alternatives are excluded.

\subsubsection*{The small-prime candidate and penalty conditions}

The source's lower transition scale is $p_0=k^2x^{26}$. Since
$k\ge5$ and $x>1$, $p_0>2x$. Its upper candidate scale is
$X'=\exp(x^{\epsilon/3})>2x$ for all sufficiently large $x$
with fixed $\epsilon>0$. Therefore, for an odd unramified prime
in \texttt{inertClass} and in the stated range, all the clauses
of \texttt{Cand} hold. In particular its conditional band clause
$p_0\le p\Rightarrow p\notin\mathrm{InBand}$ is vacuous.

The penalty condition also follows directly from the definitions.
Write $h_p$ for the $A$-coordinate. Because the prime is in
\texttt{inertClass}, the orbit step is $h_p$. The source chooses
one representative of each orbit using a map $b$ satisfying
$b(t)+j h_p=t$ for some $j$, and $b(t+h_p)=b(t)$.
These identities imply $b(b(t))=b(t)$ and that $b$ is constant
on each orbit. Distinct representatives therefore give disjoint
orbits. If $N_t$ counts top labels in the orbit of $t$, then
\[
               \sum_{t\in\mathrm{goodReps}}N_t\le |TT|.
                                                        \tag{2.8}
\]
Equivalently, the positions $0\le j<\operatorname{ord}(h_p)$
with $t+jh_p\in TT$ are in bijection with the top labels whose
representative is $t$. This verifies the equality of the two
top-count definitions used by the source, rather than assuming it.

Put $R_\epsilon=\lceil8(4+8\epsilon)/\epsilon\rceil+1$,
the source's constant \texttt{rG}, and let $c_{\mathrm{II}}$
be its fixed local constant \texttt{kII}. Below $p_0$ the
definition of the orbit penalty, independently of either word,
is
\[
 \kappa_t=\min(1,c_{\mathrm{II}}N_tR_\epsilon^2k x^{-26}),
 \qquad \Sigma=m\sum_{t\in\mathrm{goodReps}}\kappa_t.
                                                        \tag{2.9}
\]
Using (2.8), $m\le2n\le k$, and $|TT|\le n\le k$, gives
\begin{align}
 \Sigma&\le\max(0,c_{\mathrm{II}})R_\epsilon^2 m k|TT|x^{-26}
       \le\max(0,c_{\mathrm{II}})R_\epsilon^2 k^3x^{-26}\notag\\
 &\ll_{\epsilon,c_{\mathrm{II}}}
       x^{-25/2+(3003/100)\epsilon}\le x^{-\epsilon}
                                                        \tag{2.10}
\end{align}
for all sufficiently large $x$, by the upper \texttt{SizeK}
bound and $\epsilon<1/52$. No positivity of
$c_{\mathrm{II}}$ is needed for this upper bound.
The definition of \texttt{VV} now makes it exactly
\texttt{goodReps}, for every pair of word functions.

For completeness, the upper support scale is
$X''=xn\exp(x^{\epsilon/3})>2x$; here the lower degree window
gives the integer $n\ge2$. In the source definition
\texttt{NoGain}, the low-order alternative is excluded by
(2.7), the band alternative by $p<p_0$, the excessive-penalty
alternative by (2.10), the large-prime alternative by $p<X''$,
the prime-class alternative by \texttt{inertClass}, and the
odd-order alternative by \texttt{HEv}. Thus \texttt{NoGain}
is false. This proves the remaining assertions of the opening
theorem. It establishes neither a nonempty good-orbit set nor
a lower bound for the gain, and does not supply a cubic
character weighting.

The literal \texttt{FamCore} quantities \texttt{gainTot} and
\texttt{ioTot} sum local terms by \emph{rational residue
characteristic}. They do not select prime ideals of absolute norm
at most $2x$. The local data identify an orbit length with the
order of the $A$-coordinate of Frobenius, and the growth argument
explicitly uses orbits of order at least $x^{\epsilon/4}$.
These local quantities may therefore be nonzero when every ordinary
Hecke coefficient considered above vanishes. Theorem~\ref{thm:empty-norms}
rules out the ordinary low-norm comparison, including its signed
mixtures; it does not rule out that rational-prime orbit statistic.

At a fixed scale \texttt{FamCore} holds the frame, the Artin map,
the root discriminant, the auxiliary level primes, and the labelled
top set fixed across its family members. Varying the family index
does not create a different low-norm Euler support. Its covering
condition varies weight and translation data. A new comparison
must therefore identify the local orbit gain, with its
normalization, rather than substitute an ordinary Hecke sum for it.

These conclusions are ordinary mathematical consequences of the
displayed source definitions. No compilation or kernel verification
of the present implications in Lean is asserted.


\section{The numerical size conditions}

Set $\vartheta=\epsilon/100$, $L=x^\epsilon$ and
$\delta=x^{-\epsilon/8}$. We use the full numerical windows
\begin{align}
 x^{-\vartheta}\sqrt{x\delta}&\le m\le
            x^\vartheta\sqrt{x\delta},&
 \log f&\le\sqrt{m/2-1},\notag\\
 (ke^L)^8&\le g_i,&\log g_i&\le x^\vartheta L\quad(i=0,1),\notag\\
 x^{2+8\epsilon-\vartheta}&\le n\le x^{2+8\epsilon+\vartheta},&
 x^{9/2+10\epsilon-\vartheta}&\le k\le
                         x^{9/2+10\epsilon+\vartheta},\notag\\
 f&=2n,& 2\le m&\le f\le k,\qquad m+3\le k.
                                                        \tag{3.1}
\end{align}
Here $m$ is even. These are numerical consequences of the source's
\texttt{SizesAt} record; they are only part of that record. The
source's \texttt{SizeG} also asks $g_i\le2g_j$, which our choice
below satisfies. We may take $k$ even as required by its rank data.

Choose
\[
 n=2^{\lfloor\log_2(x^{2+8\epsilon})\rfloor},\quad
 m=2\left\lceil\frac{x^{1/2-\epsilon/16}}2\right\rceil,
 \quad k=2\left\lceil\frac{x^{9/2+10\epsilon}}2\right\rceil.
                                                        \tag{3.2}
\]
The factor of at most two in $n$, and the bounded rounding errors
in $m,k$, are absorbed by $x^\vartheta$ for fixed positive
$\epsilon$. Also $\log f=O(\log x)$ whereas $\sqrt m$ is a positive
power of $x$. The ordering of the exponents proves the rank
inequalities. Finally let $Y=(ke^L)^8$ and take two distinct primes
$g_0,g_1\in[Y,2Y]$. The prime number theorem supplies them for
sufficiently large $Y$, and
\[
 \log g_i\le8(L+\log k)+\log2=o(x^\vartheta L).
                                                        \tag{3.3}
\]
This proves every assertion in (3.1). No local level, lattice,
CM type, or theta lift is constructed by assigning these numbers.

It is useful to distinguish a discarded stronger target. If one
also demanded $\log[F:\Q]\asymp x^{1/4}$, then already the lower
bound $g_i\ge(ke^L)^8$, its upper logarithmic bound, and $f\le k$
would imply
\[
 \log[F:\Q]\le\log2+\tfrac18x^{101\epsilon/100}
             \le\log2+\tfrac18x^{101/5200}=o(x^{1/4}).
                                                        \tag{3.4}
\]
This excludes that additional target. It does not contradict
(3.2), whose degree has logarithm $O(\log x)$.

\section{Construction of the common CM field}

Write $n=2^{s-1}$ with $n$ as in (3.2), and put
\[
 \mathcal S_x=\{p\le2x:p\text{ prime},\ p\equiv1\pmod3\},
 \quad u=|\mathcal S_x|,
 \quad r=u+s.
                                                        \tag{4.1}
\]
The symbol $u$ here counts prime constraints and is unrelated to
the rank parameter $t$ in (1.8).
The prime number theorem in the fixed progression $1\pmod4$
gives
\[
 \#\{2x<\ell<10x:\ell\text{ prime},\ell\equiv1\pmod4\}
       \sim\frac{4x}{\log x}.
\]
Since $u\le\pi(2x)\sim2x/\log x$ and $s=O(\log x)$, there
are at least $r$ such primes for all sufficiently large $x$.
Choose distinct $\ell_1,\ldots,\ell_r$ from this interval and set
\[
 A=K(\sqrt{\ell_1},\ldots,\sqrt{\ell_r}),\qquad
 V=\Gal(A/K)\simeq\F_2^r.
                                                        \tag{4.2}
\]
The prime squareclasses are independent over $\Q$; their compositum
is totally real and has degree $2^r$. Its intersection with the
imaginary quadratic field $K$ is $\Q$. Thus
$\Gal(A/\Q)\simeq C_2^{r+1}$, which also proves (4.2).

For $p\in\mathcal S_x$, its two primes in $K$ have the same
Frobenius in $V$, because $A/\Q$ is abelian. In the usual
coordinates it is
\[
 v_p=\left(\frac{1-(\ell_j/p)}2\right)_{j=1}^r.
                                                        \tag{4.3}
\]
Their span has dimension at most $u$. Extend that span to a
subspace $H\subset V$ of dimension exactly $u$, and define
\[
                              F_x=A^H.
                                                        \tag{4.4}
\]
Every subgroup of the abelian Galois group of $A/\Q$ is normal,
so $F_x/\Q$ is abelian. If $A^+=\Q(\sqrt{\ell_1},\ldots,
\sqrt{\ell_r})$, then $F_x=K(A^+)^H$ and
\[
 [F_x:K]=2^s=h,\qquad [F_x:\Q]=2^{s+1}=4n.
                                                        \tag{4.5}
\]
The real field $(A^+)^H$ has degree $h$, so $F_x$ is a totally
imaginary quadratic extension of a totally real field: it is CM.

Every $p\in\mathcal S_x$ has Frobenius in $H$ and therefore
splits completely in $F_x$. All rational primes at most $2x$
other than $3$ are unramified in $F_x$, since its only possible
ramification is at $3$ and the $\ell_j$. In particular $2$ is
unramified: the discriminant of each real quadratic generator
is $\ell_j$, because $\ell_j\equiv1\pmod4$.
If $p\equiv2\pmod3$, its Frobenius in $\Gal(F_x/\Q)$ is
nontrivial on $K$ and hence nonidentity. Its order is exactly
two because the whole group has exponent two. Its relative
Frobenius over the inert prime of $K$ is its square and is
identity. This proves assertion 2 of Theorem~\ref{thm:construction}.

\subsection{Exact discriminant and cubic conductor}

Let $W=H^\perp\subset\F_2^r$, so $\dim W=s$. For $w\in W$ put
\[
 D_w=\prod_{j=1}^r\ell_j^{w_j},\qquad D_0=1,
 \quad J=\{j:\text{some }w\in W\text{ has }w_j=1\}.
                                                        \tag{4.6}
\]
The characters of $F_x/\Q$ are exactly the real quadratic
characters of conductor $D_w$ and their products with
$\eta=\chi_{-3}$, of conductor $3D_w$. These conductor assertions
follow from the quadratic discriminant formula, since every
$D_w\equiv1\pmod4$ and is coprime to $3$.
The conductor--discriminant formula
\citep[Chapter V, Theorem~3.27]{Milne2020} gives
\[
 |D_{F_x}|=3^h\prod_{w\in W}D_w^2
           =3^h\prod_{j\in J}\ell_j^h.
                                                        \tag{4.7}
\]
For the second equality, each nonzero coordinate functional on
$W$ is one on exactly half of its $h$ elements. Taking the
$2h$-th root proves (1.5), with
\[
 \log\rd(F_x)\le\tfrac12\log3+\tfrac r2\log(10x)=O(x).
                                                        \tag{4.8}
\]
The conductor of $F_x/\Q$ divides $3\prod_j\ell_j$.

The polynomial $X^3-269$ is Eisenstein at $269$, and its
discriminant $-27\cdot269^2$ is not a square. Thus its splitting
group is $S_3$, whose maximal abelian subfield is $K$.
For any abelian $F/\Q$ containing $K$, the intersection $E\cap F$
is therefore $K$. In particular $EF/F$ is cyclic of degree three
and $\chi_F$ is faithful.

The extension $E/K$ is unramified away from $3,269$. It splits
at $3$: for $f(X)=X^3-269$, the approximation $-1$ satisfies
$v_3(f(-1))=3>2v_3(f'(-1))=2$, so Hensel's lemma gives a root
in $\Q_3$. All three roots then lie in $K\otimes\Q_3$.
At $269$ it is totally and tamely ramified of degree three,
so a faithful character has conductor exponent one.
For $x>269$, the field $F_x$ is unramified at $269$.
Unramified base change preserves this tame degree-three inertia,
while base change creates no ramification at any other prime
for this character. This proves (1.3).
The arithmetic conductor of the induced $L$-function is exactly
\[
       Q(\chi_{F_x})=|D_{F_x}|\,269^{2h},\qquad
 \frac{\log Q(\chi_{F_x})}{h}
          =2\log\rd(F_x)+2\log269.
                                                        \tag{4.9}
\]
Its archimedean factor is $\Gamma_{\mathbb C}(s)^h$, in the
convention $\Gamma_{\mathbb C}(s)=2(2\pi)^{-s}\Gamma(s)$.

\section{Exact transfer, and what pointwise transfer forces}

For each $p\le2x$, $p\notin S$, every prime $\mathfrak p$ of $K$
above $p$ splits into $h$ primes of $F_x$, each of the same norm
and character value. Multiplication of the Euler factors gives
(1.4) directly. Differentiating their logarithms proves (1.1)
coefficientwise, for every power of each such $p$, even if the
power exceeds $2x$.

More explicitly, if $p\equiv1\pmod3$ and
$z=\chi(\mathfrak p)\in\{1,\omega,\omega^2\}$, the normalized
coefficient at $p^k$ is
\[
                         (z^k+z^{-k})\log p.
                                                        \tag{5.1}
\]
If $p\equiv2\pmod3$, its Frobenius in $S_3$ is a transposition,
whose square is identity. Thus the character of its degree-two
prime in $K$ is one. There are $h$ primes of norm $p^2$ in
$F_x$, and the normalized coefficient is $2\log p$ for even
$k$, zero for odd $k$. This is exactly $a_{p,k}\log p$.
No higher-prime-power error is needed for the full logarithmic
derivative. In particular for any complex weights $c_n$ supported
on $n\le2x$,
\[
                 \frac1h\sum_n c_n b_{F_x}(n)
                       =\sum_n c_n b_K(n).
                                                        \tag{5.2}
\]
The equality has no degree-dependent constant and no sum over hosts.

\begin{lemma}\label{lem:forced-splitting}
Let $F/\Q$ be abelian, contain $K$, and let $p\equiv1\pmod3$
be a rational prime. Write $e_p,f_p$ for its common ramification
and residue degrees in $F$. Then the left side of (1.6) is zero
if $f_p>1$, and is $a_p/e_p$ if $f_p=1$.
Consequently (1.6) forces $e_p=f_p=1$.
\end{lemma}
\begin{proof}
Since $p\notin\{3,269\}$, the character $\chi$ is unramified at
both primes of $K$ above $p$, and remains unramified after base
change. If $f_p>1$ there are no degree-one primes to sum.
If $f_p=1$, each of the two primes of $K$ has $h/e_p$ primes
above it in $F$. Their character values are respectively
$z$ and $z^{-1}$, by compatibility with the norm of prime ideals.
The normalized sum is $(z+z^{-1})/e_p=a_p/e_p$.
Since $a_p$ is $2$ or $-1$, equality with $a_p$ forces $e_p=1$.
This proof also covers a host initially allowed to ramify at $p$.
\end{proof}

\section{A uniform lower bound for the discriminant}

\subsection{A small split prime detects every extra abelian character}
\label{sec:abelian-character-detector}

The following argument is conditional on Dirichlet GRH. It does not require GRH for a nonabelian extension.
Write $\eta=\chi_{-3}$.

\begin{proposition}\label{prop:abelian-character-detector}
Assume GRH for Dirichlet $L$-functions.
There is an absolute constant $C>0$ with the following property.
Let $\xi$ be a primitive Dirichlet character of conductor $q$,
and suppose $\xi\ne1,\eta$. If
\[
 \xi(p)=1
 \qquad\text{for every rational prime }p\le 2x
       \text{ with }p\equiv1\pmod3,\ p\nmid q,
                                                        \tag{A.1}
\]
where $x\ge2$, then
\[
                         x\le C\log^2(3q).              \tag{A.2}
\]
In particular, $\log(3q)\gg\sqrt{x}$.
The same conclusion holds if the condition in (A.1) is imposed
only for the primes in $[x,2x]$.
\end{proposition}

\paragraph{Proof.}
A nonnegative prime coefficient detects precisely a failure of
$\xi(p)=1$ at a prime split in $\mathbb Q(\sqrt{-3})$.
Its smooth sum has a positive pole term of order $x$.
Dirichlet GRH bounds its zero terms by
$O(\sqrt{x}\log(3q))$, uniformly in the conductor.
Ramified primes and higher powers do not have order $x$.
We supply the uniform zero estimate, including the imprimitive
product $\eta\xi$.

Fix once and for all a nonzero $W\in C_c^\infty((1,2))$ with
$W\ge0$, and write
\[
 \widehat W(s)=\int_0^\infty W(u)u^{s-1}\,du,\qquad
                         w_1=\widehat W(1)>0.
\]
Its Mellin transform decreases faster than every power of $|\Im s|$
on each fixed vertical strip.

We first prove that for a primitive Dirichlet character $\chi$
of conductor $d$, including the character of conductor one,
\[
 \sum_{n\ge1}\Lambda(n)\chi(n)W(n/x)
 ={\bf1}_{\chi=1}w_1x+
             O_W\bigl(\sqrt{x}\log(3d)\bigr).
                                                        \tag{A.3}
\]
For $\chi\ne1$, set $a\in\{0,1\}$ by $\chi(-1)=(-1)^a$ and
\[
 \Lambda_\chi(s)=
 (d/\pi)^{(s+a)/2}\Gamma((s+a)/2)L(s,\chi).
\]
This is entire of order one and satisfies
$\Lambda_\chi(s)=\epsilon_\chi
 \overline{\Lambda_\chi(1-\overline s)}$, with
$|\epsilon_\chi|=1$.
For $\chi=1$ use instead
$\xi_{\mathbb Q}(s)=s(s-1)\pi^{-s/2}\Gamma(s/2)\zeta(s)$.
These analytic properties follow from the Dirichlet case of
Tate's functional equation; see \citet[Theorems~5.16 and 5.22(b),(c), pp.~34--39]{Poonen2015}.

Here is the required local zero count, with its conductor dependence.
The Hadamard product and the functional equation give
\[
 \sum_\rho
 \frac{2-\Re\rho}{(2-\Re\rho)^2+(t-\Im\rho)^2}
 =
 \frac12\log(d/\pi)
 +\frac12\Re\frac{\Gamma'}{\Gamma}\left(\frac{2+a+it}{2}\right)
 +\Re\frac{L'}{L}(2+it,\chi)                              \tag{A.4}
\]
for $\chi\ne1$. Zeros are counted with multiplicity.
To see that no unknown constant occurs, differentiate the
genus-one Hadamard product. Its real part is the displayed zero
sum plus a real constant. On $\Re s=1/2$ the functional equation
makes the logarithmic derivative purely imaginary. The zero
terms there cancel in the pairs $\rho,1-\overline\rho$;
under the granted GRH each individual real part is zero.
The constant is therefore zero. The sums of real parts converge
absolutely, since the function has order one and $0<\Re\rho<1$.
For $\xi_{\mathbb Q}$ the identity has the additional terms
$\Re(1/(2+it)+1/(1+it))$, which are bounded.

At $\Re s=2$, the logarithmic derivative of the Euler product
is bounded absolutely by $\sum_n\Lambda(n)n^{-2}$.
The gamma logarithmic derivative is
$O(\log(3+|t|))$ uniformly on the indicated line.
Every zero with $|\Im\rho-t|\le1$ contributes at least $1/5$
to the left side of (A.4). Consequently
\[
 \#\{\rho:|\Im\rho-t|\le1\}
       \ll \log\bigl(3d(3+|t|)\bigr).                    \tag{A.5}
\]
All constants here are absolute.
This also proves
\[
 \sum_\rho|\widehat W(1/2+i\Im\rho)|
                        \ll_W \log(3d),                 \tag{A.6}
\]
by summing (A.5) over unit height intervals and using Mellin
decay of order, for example, four.

Mellin inversion of $-L'/L$ on $\Re s=2$ and contour displacement
to $\Re s=-1/2$ give
\[
 \sum_{n\ge1}\Lambda(n)\chi(n)W(n/x)
 ={\bf1}_{\chi=1}w_1x
   -\sum_\rho x^\rho\widehat W(\rho)+O_W(\log(3d)).
                                                        \tag{A.7}
\]
For completeness, the displaced integral is
$O_W(x^{-1/2}\log(3d))$: the functional equation expresses
$L'/L(-1/2+it,\chi)$ using the absolutely convergent Euler
product at $3/2-it$ and the gamma logarithmic derivatives.
The possible trivial zero at zero contributes $O_W(1)$.
Choose horizontal heights in successive unit intervals a distance
$\gg1/\log(3d(3+|t|))$ from all zero ordinates, using (A.5).
The logarithmic derivative obtained from the Hadamard product
has polynomial logarithmic growth on those contours, and the
rapid Mellin decay makes their integrals tend to zero.
This justifies the displacement and absolute convergence of its
zero sum. Under GRH, $|x^\rho|=\sqrt{x}$.
Equations (A.6) and (A.7) imply (A.3).

If a character modulo $m$ is induced from a primitive character
of conductor $d\mid m$, replacing it by that primitive character
changes its smooth prime-power sum only at $p\mid m$.
For a fixed $p$, the number of exponents $k$ with
$x<p^k<2x$ is at most $1+\log2/\log p$. Hence the total change is
\[
 \ll_W \sum_{p\mid m}\log p
 \ll_W\log(3m).                                         \tag{A.8}
\]
Thus (A.3), with $d$ replaced by $m$, also holds for an
imprimitive character, retaining its correct principal pole term.

Define, on prime powers,
\[
 h(p^k)=\frac{1+\eta(p)^k-\Re(\xi(p)^k)
                    -\Re((\eta\xi)(p)^k)}2,
\]
where the usual zero extension is used at primes dividing
a character's modulus. Equivalently use the four Dirichlet
characters $1,\eta,\xi,\eta\xi$ at the integer $p^k$.
The last three are nonprincipal: the assumptions
$\xi\ne1,\eta$ exclude both possible principal products.
The primitive conductor inducing $\eta\xi$ divides $3q$.
Applying (A.3) and (A.8) therefore gives
\[
 \sum_{p,k\ge1}(\log p)h(p^k)W(p^k/x)
       =\frac{w_1x}{2}
                       +O_W(\sqrt{x}\log(3q)).           \tag{A.9}
\]
For $p\nmid3q$,
\[
 h(p)=\frac{(1+\eta(p))(1-\Re\xi(p))}{2}\ge0.
\]
It is zero at inert primes and, by (A.1), at every split
prime in the support of $W(p/x)$. The exceptional prime terms
with $p\mid3q$ contribute $O_W(\log(3q))$.

Finally $|h(p^k)|\le2$, and the elementary Chebyshev bound
$\vartheta(y)=\sum_{p\le y}\log p\ll y$ yields
\[
 \sum_{\substack{k\ge2\\p^k\le2x}}\log p
 \le\sum_{2\le k\le\log_2(2x)}\vartheta((2x)^{1/k})
 \ll \sqrt{x}+x^{1/3}\log(3x)
 \ll\sqrt{x}.                                           \tag{A.10}
\]
The last inequality holds with an absolute constant for
all $x\ge2$. Inserting these bounds into (A.9) gives
$w_1x/2\ll_W\sqrt{x}\log(3q)$, proving (A.2).

\paragraph{Conductor--discriminant consequence.}
If $F/\mathbb Q$ is abelian and contains
$K=\mathbb Q(\sqrt{-3})$, and every degree-one prime of $K$
above a rational prime $p\le2x$ splits completely in $F/K$,
then every character in the character group of $F/\mathbb Q$
other than $1,\eta$ has conductor $q_\xi$ satisfying
$\log(3q_\xi)\gg\sqrt{x}$.
By the conductor--discriminant formula,
\[
 \log\operatorname{rd}(F)
   =\frac1{[F:\mathbb Q]}\sum_\xi\log q_\xi.
\]
For $F\ne K$, its degree is an even integer at least four,
so at least half its characters are outside $\{1,\eta\}$.
It follows, for all sufficiently large $x$, that
$\log\operatorname{rd}(F)\gg\sqrt{x}$, with an absolute
constant. This last conclusion needs the stated complete
splitting condition; preserving a selected trace through
some other mechanism is not asserted to imply that condition.



\begin{proof}[Proof of Theorem~\ref{thm:discriminant}]
Lemma~\ref{lem:forced-splitting} makes every prime in question
split completely in $F$. Every Dirichlet character attached to
$F/\Q$ is then one at these primes. Apply
Proposition~\ref{prop:abelian-character-detector} to each primitive
character other than $1,\eta$. If $D=[F:\Q]\ge4$, at least
$D-2\ge D/2$ characters have this property. For an absolute
$c_0>0$, their conductors satisfy
$\log\cond(\xi)\ge c_0\sqrt{x}-\log3$.
The conductor--discriminant formula gives
\[
 \log\rd(F)=\frac1D\sum_\xi\log\cond(\xi)
 \ge\left(1-\frac2D\right)(c_0\sqrt{x}-\log3)
                       +\frac{\log3}{D}.
                                                        \tag{6.1}
\]
Increasing an absolute threshold makes this at least
$c_0\sqrt{x}/4$. This proves the uniform assertion.
\end{proof}

\subsection{An explicit bound for the multiquadratic construction}

For this construction the bound can be made explicit directly from
\citet[Theorem~5.1, p.~1729]{BachSorenson1996}, without the general
character argument. If $0\ne w\in W$, the biquadratic field
$B_w=K(\sqrt{D_w})$ has degree four and discriminant
$9D_w^2$. There is an automorphism of $B_w$ trivial on $K$ and
nontrivial on $\sqrt{D_w}$. Under Dirichlet GRH, Bach and
Sorenson give an unramified rational prime with this Frobenius and
\[
 p\le(4\log|D_{B_w}|+2.5\,[B_w:\Q]+5)^2
       =(8\log D_w+4\log9+15)^2.
                                                        \tag{6.2}
\]
The hypothesis of their theorem is satisfied because the zeta
function of $B_w$ is a product of Dirichlet $L$-functions.
Such a prime is $1\pmod3$ but is not split in $F_x$, so $p>2x$.
It follows that
\[
 \log D_w>\frac{\sqrt{2x}-4\log9-15}{8}.
                                                        \tag{6.3}
\]
Averaging (6.3) in (4.7) yields the explicit inequality
\[
 \log\rd(F_x)>\frac12\log3+
       \left(1-\frac1h\right)
       \frac{\sqrt{2x}-4\log9-15}{8}.
                                                        \tag{6.4}
\]
This is a conditional lower bound, while the construction and
upper bound (4.8) are unconditional.

\section{The literal discriminant and group conditions}

We record precisely the mathematical content of the source
conditions used here. In the formalism of \citet{AlpogeLean2026},
the host is a number field with an embedding into a cyclotomic
field of conductor $Q_Sqd$. In particular it is abelian over $\Q$.
Its root discriminant parameter is defined to be
\[
       R_E=|\operatorname{disc}(E_{\mathrm{host}})|^%
                              {1/[E_{\mathrm{host}}:\Q]}.
                                                        \tag{7.1}
\]
It is not a free height parameter. The source's $q,d$ are primes
and $Q_S$ is a positive integer. Its scale definitions are
$L=x^\epsilon$, $b=\exp L$. The clauses called
\texttt{SizeT}, \texttt{SizeCond}, and \texttt{SizeQD} include
exactly
\[
 t\ge1,\qquad (Q_Sqd)^{36}t^{150}\le\sqrt b,
 \qquad 1\le R_E\le\frac{Q_S}{2}\sqrt{qd}.
                                                        \tag{7.2}
\]
Only these displayed parts are needed; other clauses cannot
weaken a necessary consequence of them.

\begin{lemma}[Literal upper bound]\label{lem:upper}
The inequalities (1.8), and hence (7.2), imply
\[
                  \log\rd(F)\le\frac{x^\epsilon}{72}-\log2.
                                                        \tag{7.3}
\]
\end{lemma}
\begin{proof}
Since $q,d\ge1$, one has $2\rd(F)\le Q\sqrt{qd}\le Qqd$.
Since $t\ge1$, the conductor condition gives
$(2\rd(F))^{36}\le e^{x^\epsilon/2}$. All quantities are
positive, so taking logarithms and dividing by $36$ proves (7.3).
\end{proof}

The group condition is separate. A source frame supplies a finite
abelian group $A$, a natural number $e$, and a surjection
\[
 A\longrightarrow
 H_e=(\Z/2^{e+2}\Z)^\times/\{\pm1\}\simeq C_{2^e}.
                                                        \tag{7.4}
\]
The host data require $e\ge2$ and an isomorphism
\[
            \Gal(E_{\mathrm{host}}/\Q)\simeq C_2^2\times A.
                                                        \tag{7.5}
\]
The elementary structure of the units modulo $2^{e+2}$ proves
the cyclic description in (7.4): every odd class is uniquely
$\pm5^j$ modulo $2^{e+2}$, with $j$ modulo $2^e$.
Composing the projection in (7.5), the surjection in (7.4), and
$C_{2^e}\twoheadrightarrow C_4$ shows that the literal host has
a cyclic quotient of order four. A group of exponent two has
no such quotient. This proves the unconditional structural part
of Theorem~\ref{thm:literal}; it also shows that a literal host
has degree greater than two.

\begin{proof}[Proof of the remaining part of Theorem~\ref{thm:literal}]
The literal host is abelian over $\Q$, has degree greater than
two, and, by assumption, contains $K$ and satisfies (1.6).
Theorem~\ref{thm:discriminant} gives
$\log\rd(F)\ge c\sqrt{x}$, while Lemma~\ref{lem:upper} gives
the upper bound (7.3). They are incompatible for all sufficiently
large $x$, because $\epsilon<1/52<1/2$.
The same proof applies to the more general scalar formulation.
\end{proof}

No comparison such as $g_i\ge\rd(F)$ has been assumed or is
needed. In the source, the $g_i$ are two distinct auxiliary level
primes. The discriminant restriction comes instead from (7.2).
The argument is an ordinary mathematical inspection of the
specified definitions. It is not a claim that this implication
has been compiled or kernel-checked in Lean.

\section{The original one-prime cyclotomic construction}

For completeness, another natural candidate has an exact but
different limitation. Let $q\equiv1\pmod3$ be a rational prime,
$M\ge2$, $N=q^M-1$, and
\[
             F_{q,M}=\Q(\zeta_N)^{\langle q\bmod N\rangle}.
                                                        \tag{8.1}
\]
The order of $q\bmod N$ is exactly $M$: for $0<j<M$,
$0<q^j-1<N$. Its subgroup has least positive representatives
$1,q,\ldots,q^{M-1}$ and does not contain $-1=N-1$.
Thus $F_{q,M}$ is CM, abelian over $\Q$, and has degree
$\varphi(N)/M$. Since $3\mid N$ and $q\equiv1\pmod3$, it
contains $K$. The prime $q$ is unramified and its Frobenius is
trivial in the quotient, so it splits completely in $F_{q,M}$.
As before $E\cap F_{q,M}=K$ and the cubic character stays faithful.

If $q\in[x,2x]$ and $x$ is sufficiently large, then $N>2x$.
For an unramified prime $p\le2x$, complete splitting is equivalent
to $p\bmod N\in\langle q\rangle$. Among
$1,q,\ldots,q^{M-1}$, the only rational prime is $q$ itself.
Thus this field preserves exactly one degree-one rational prime
in the target range. Combining one such field for each $q$
requires summing the separate estimates. A uniform bound $B$
for each normalized sum gives only $|\mathcal S_x|B$ by the
triangle inequality. No cancellation among that collection is
asserted here. The common field (4.4) avoids this counting loss.

There is also a degree-window failure for (8.1). If $M=2$, then
$6\mid N$ and $\varphi(N)\le N/3$, so
\[
 [F_{q,2}:\Q]\le(q^2-1)/6<\tfrac23x^2.
                                                        \tag{8.2}
\]
If $M\ge3$, the elementary bound $\varphi(N)\ge N^{3/4}/3$
gives
\[
 [F_{q,M}:\Q]\ge\frac{q^{3M/4}}{3\,2^{3/4}M}
             \ge\frac{x^{9/4}}{9\,2^{3/4}}.
                                                        \tag{8.3}
\]
To prove the totient bound, for every prime $p\ge5$ one has
$1-1/p\ge p^{-1/4}$; the factors at $2,3$ have product at least
$1/3$. Multiplying gives the result, including higher prime
powers. The last inequality in (8.3) follows because
$q^{3M/4}/M$ increases for $M\ge3$ and $q\ge7$.
The upper degree window in (3.1) is
$[F:\Q]\le4x^{2+801\epsilon/100}$, whose exponent is less
than $2+801/5200<9/4$. Thus neither case meets the full
degree window for large $x$. A weakened set of numerical assumptions that
retains only the lower bound on $n$ does not give this conclusion.

The exact conductor can also be stated when $F_{q,M}$ ramifies
at $269$. If $e_{269}$ is its common ramification index there,
then
\[
 \mathfrak f(\chi_{F_{q,M}})
       =\prod_{\mathfrak P\mid269}\mathfrak P,\qquad
 \N\mathfrak f=269^{[F_{q,M}:\Q]/e_{269}}.
                                                        \tag{8.4}
\]
Indeed any abelian ramification index at $269$ divides a number
$268\cdot269^a$, prime to three, so tame cubic inertia survives
base change. If $N=269^a R$, $(269,R)=1$ and $a\ge1$, then
\[
 e_{269}=\frac{\varphi(269^a)\operatorname{ord}_R(q)}{M},
                                                        \tag{8.5}
\]
with $\operatorname{ord}_1(q)=1$; if $a=0$ it is one.
This follows by intersecting $\langle q\rangle$ with the
$269$-inertia subgroup in the cyclotomic Galois group.

\section{Weighted mixtures under relaxed host conditions}
\label{sec:weighted}

For an abelian $F\supset K$, define
\[
 \alpha_F(p)=
 \begin{cases}e_p^{-1},&f_p=1,\\0,&f_p>1,
 \end{cases}
 \qquad (p\equiv1\pmod3).
                                                        \tag{9.1}
\]
Lemma~\ref{lem:forced-splitting} says that its normalized
degree-one cubic trace is $a_p\alpha_F(p)$. Thus $0\le\alpha_F\le1$,
including primes ramified in $F$.

\begin{proposition}[Positive mixtures]
If $c_i\ge0$, $\sum_i c_i=1$, and
$\sum_i c_i\alpha_{F_i}(p)=1$ on a set of split-$K$ primes,
then every field with $c_i>0$ splits completely at every prime
in that set.
\end{proposition}
\begin{proof}
For each prime, $\sum_i c_i(1-\alpha_{F_i}(p))=0$ is a sum of
nonnegative numbers. Each active summand is zero, so
$\alpha_{F_i}(p)=1$, which means $e_p=f_p=1$.
The same proof applies to a countable nonnegative mixture.
\end{proof}

For literal hosts Theorem~\ref{thm:empty-norms} is stronger and
does not require positivity. The next estimate remains useful
after the common cyclotomic condition is removed.

\begin{proof}[Proof of Theorem~\ref{thm:mixture-cost}]
Let $D=[F:\Q]$ and write $\Lambda_F(n)$ for the coefficients of
$-\zeta_F'/\zeta_F$. Abelian factorization expresses $\zeta_F$
as the product of the primitive Dirichlet $L$-functions of its
$D$ characters. Apply the uniform smooth explicit formula proved
in Section~6 to each character, and divide by $D$. The single
principal character supplies the main term. The
conductor--discriminant formula gives
\begin{align}
 \frac1D\sum_n\Lambda_F(n)W(n/x)
 &=\frac{xI}{D}
       +O_W\!\left(\sqrt{x}\left(1+
                    \frac1D\sum_\xi\log\cond(\xi)\right)\right)\notag\\
 &=\frac{xI}{D}+O_W\bigl(\sqrt{x}(1+\log\rd(F))\bigr),
 \qquad I=\int_1^2W(u)\,du>0.
                                                        \tag{9.2}
\end{align}
The implied constant is independent of $D$, $F$, and the number
of characters. Complex characters cause no problem: the estimates
are absolute and their sum is the real expression on the left.

Every $\Lambda_F(n)$ is nonnegative. At a split-$K$ prime $p$,
its coefficient at $n=p$, divided by $D$, is exactly
$\alpha_F(p)\log p$: if $f_p=1$, there are $D/e_p$ prime ideals
of norm $p$; otherwise there are none. Discarding all other
terms in (9.2) yields
\[
 \sum_{p\equiv1(3)}\alpha_F(p)\log p\,W(p/x)
 \le\frac{xI}{D}+C_W\sqrt{x}(1+\log\rd(F)).
                                                        \tag{9.3}
\]
With $D\ge c x^2$ and $\log\rd(F)\le x^\epsilon/72$, this is
at most $C_{W,c}x^{1/2+\epsilon}$ for $x\ge1$.

Put $A(x)=\sum_{p\equiv1(3)}\log p\,W(p/x)$.
The prime number theorem in this fixed progression gives
$A(x)\sim Ix/2$. For an exact mixture, the triangle inequality
and (9.3) give
\[
 A(x)\le\sum_i|c_i|
       \sum_{p\equiv1(3)}\alpha_{F_i}(p)\log p\,W(p/x)
       \le C_{W,c}x^{1/2+\epsilon}\sum_i|c_i|.
                                                        \tag{9.4}
\]
This proves the claimed lower bound. If the weighted absolute
defect is at most $A(x)/4$, the first left side in (9.4) may
instead be $3A(x)/4$, and the conclusion is unchanged.
\end{proof}

This is a cost bound for pointwise recovery or weighted absolute
approximation. Applying only separate absolute estimates to the
host terms incurs this coefficient mass. It does not assert that
one particular signed aggregate cannot cancel or admit a different
comparison. The unconditional zero statement for literal hosts
has the stronger, but precisely ordinary-Euler-product, scope of
Theorem~\ref{thm:empty-norms}.

\section{A fixed quotient and its exact Artin projector}
\label{sec:projector}

Put $F_0=K\Q(\zeta_{16})^+$ and
$G_0=\Gal(F_0/K)\simeq C_4$. The factors in
Theorem~\ref{thm:projector} are finite-order Hecke $L$-functions.
Their entireness and functional equations follow from Tate's
theory; see \citet[Theorem~5.22]{Poonen2015}. No factor is principal,
because $\chi$ has order three and $\psi$ has order dividing four.

\begin{proof}[Proof of Theorem~\ref{thm:projector}]
Let $G=\Gal(F/K)$, $H=\Gal(F/F_0)$, and $r=|H|$. Since
$\chi$ is already a character of $G_K$, the tensor identity for
induction gives
\[
 V=\operatorname{Ind}_{G_F}^{G_K}(\chi|_{G_F})
        \simeq\chi\otimes\mathbb C[G].
                                                        \tag{10.1}
\]
For clarity, the isomorphism is the usual coset model: after
multiplying the component at a representative $g$ by $\chi(g)$,
the induced action becomes the scalar $\chi(\gamma)$ times
left translation by the image of $\gamma$ in $G$.
Right translations commute with this action. Averaging them over
$H$ gives the orthogonal idempotent
\[
          P_H=\frac1r\sum_{h\in H}R_h.
                                                        \tag{10.2}
\]
Its image is $\chi\otimes\mathbb C[G/H]$, of dimension four.
The regular representation of the abelian quotient is the direct
sum of its four characters. Taking the determinant of
$1-\operatorname{Frob}_{\mathfrak p}T$ on inertia invariants,
prime by prime, proves
\[
 J_0(s)=L_{F_0}(s,\chi_{F_0})
        =\prod_{\psi\in\widehat G_0}L_K(s,\chi\psi).
                                                        \tag{10.3}
\]
The same argument proves the asserted trace identity at every
prime power. At a ramified prime use $V^{I_{\mathfrak p}}$;
$P_H$ commutes with inertia, and
$(\operatorname{im}P_H)^{I_{\mathfrak p}}
=P_H(V^{I_{\mathfrak p}})$.

For the integral assertion, let $\mathcal O=\Z[\omega]$.
For each coset $C$ of $H$ in $G$, put $v_C=\sum_{g\in C}e_g$.
These four vectors are a basis of $\mathcal O[G]^H$, and
$v_C/r$ are a basis of $P_H\mathcal O[G]$. The quotient is
therefore $(\mathcal O/r\mathcal O)^4$; its determinant ideal
is $(r)^4$ and its cardinality is $r^8$. This proves the exact
index, including the coefficient-ring normalization. Each $v_C$
has Euclidean norm $\sqrt r$, and $P_H$ has norm one.

Finally $L_K(s,\chi)$ occurs as the factor $\psi=1$, and every
other factor is entire, so its zeros cannot be cancelled in
$J_0$. Thus RH for $J_0$ implies RH for the target. The dual
collection equals the original collection: conjugating ideals
of $K$ changes $\chi^{-1}$ into $\chi$, while the characters
coming from the abelian field $\Q(\zeta_{16})^+$ are unchanged
under this conjugation and are permuted by inversion. Its
functional equation therefore also shows that zero-freeness
for $\Re s>1/2$ suffices for RH.
\end{proof}

At an unramified split-$K$ rational prime the first coefficient
of $-J_0'/J_0$ is $4a_p\log p$ if $p\equiv\pm1\pmod{16}$,
and is zero otherwise. The potentially visible split classes
are exactly $1,31$ modulo $48$. Hence recovering all eight
split classes individually is stronger than necessary for a
fixed product that contains the target factor.

\subsection{An explicit sufficient analytic comparison}

This fixed product gives a precise possible replacement for an
ordinary host transfer. Let
\[
 -J_0'(s)/J_0(s)=\sum_{n\ge1}c(n)n^{-s},\qquad
 S_0(x)=\sum_{n\ge1}c(n)e^{-n/x}.
                                                        \tag{10.4}
\]
The finite-image degree-eight Artin representation gives
$|c(n)|\le8\Lambda(n)$, including at ramified prime powers.
If a construction proved $S_0(x)=O_\eta(x^{1/2+\eta})$ for every
$\eta>0$, it would prove RH for the target. Indeed, for $\Re s>1$,
absolute convergence and the substitution $u=n/x$ give
\[
 \int_0^\infty S_0(x)x^{-s-1}\,dx
                       =-\Gamma(s)J_0'(s)/J_0(s).
                                                        \tag{10.5}
\]
The proposed bound, together with exponential decay as $x\to0$,
makes the left side holomorphic on $\Re s>1/2$. Since
$\Gamma$ has no zeros or poles there, the logarithmic derivative
has no poles, and the entire function $J_0$ has no zeros there.
Its functional equation completes the assertion. The bound in
this paragraph is a sufficient condition, not an estimate proved
by the present paper or by the abstract projector.

\subsection{Normalization, denominators, and gain}

The coefficients in (10.2) have absolute sum one relative to the
\emph{unnormalized} twisted traces
$\tr(R_h\operatorname{Frob}^j\mid V)$. If these traces have
already been divided by $|G|$, recovering (10.3) instead has
coefficient mass $|G|$. Norm one is not a free normalization
of an estimate on normalized host traces.

For a general $H$-stable torsion-free $\Z$-lattice $M$ with
$a=\operatorname{rank}_{\Z}M^H$, averaging gives
\[
 M^H\subset P_HM\subset r^{-1}M^H,
 \qquad [P_HM:M^H]\le r^a.
                                                        \tag{10.6}
\]
The inclusions follow directly by multiplying the average by
$r$. Thus the denominator cost is at most $a\log r$, not
unconditionally $\log r$ for a large theta module. In the
rank-four regular example the logarithmic cost divided by
$[\Q(\omega):\Q]=2$ is exactly $4\log r$.

Orthogonal projection preserves nonnegativity of the trace of
a commuting positive operator. It does not make Frobenius
positive. A determinant bound alone does not force its divisibility
to occur on the invariant component. For an exact counterexample
take $M=\Z e_+\oplus\Z e_-$ with a $C_2$-action of signs
$+1,-1$. The equivariant integral map
\[
                 T=\operatorname{diag}(1,p^N)
                                                        \tag{10.7}
\]
has full cokernel of order $p^N$ and invariant-component
cokernel of order one, for every $N\ge1$.

Uniform lattice inclusion is different and gives a constructive
inheritance statement. If $f:M_{\mathbb Q}\to N_{\mathbb Q}$,
projectors $P,Q$ satisfy $fP=Qf$, and
$f(M)\subset\pi^a N$ over a discrete valuation ring, then
\[
                    f(PM)=Qf(M)\subset\pi^a QN.
                                                        \tag{10.8}
\]
Thus the same valuation gain survives on the projected lattices.
There is no new lower-gain theorem to prove for such an inclusion
once the coherent maps and lattices have been constructed. Comparing
$PM,QN$ with invariant integral lattices may still incur the
denominator cost in (10.6).

This distinction applies to the source. Its \texttt{vGeLat}
asserts that a vector lies in
$\pi^{\lceil g v_p(p)\rceil}N$, and \texttt{LatInst.vKge}
applies this to the image of the chain map. In a discrete valuation
ring that scaled lattice is independent of the chosen uniformizer.
A statement for every vector in a given source lattice is therefore
a uniform inclusion of the form above, rather than merely a
cokernel determinant. The unresolved task is to construct compatible
projectors on those source and target lattices, preserve the
relevant source restriction, and identify the required trace and
normalization. If the known inclusion applies on that source,
its valuation automatically passes to the projected lattices.

\section{The highest weights prevent the naive deck action}
\label{sec:highest}

We now prove the specific obstruction in
Theorem~\ref{thm:highest}. All assertions concern one fixed
family with scale-independent $C_c$. The source has
$G=\Gal(F/\Q)$ of order $4n$ and a CM-label subgroup
$\Psi\subset G$ of index two. Its top set is
$\mathcal T=\{\tau_1(t):t\in TT\}\subset\Psi$, with
$\tau_1$ injective. The source size and weight conditions give
\begin{align}
 0<x^{-\epsilon/100}n\delta
     &\le |\mathcal T|\le2n\delta,
       &\delta&=x^{-\epsilon/8},\notag\\
 a_T(t)&\ge a=x\exp(x^\epsilon),
       &S_c&=4k\Delta,\notag\\
 kS_c&\le w''(n),
       &w''(n)&=\frac{a}{xn\exp(x^{\epsilon/3})}.
                                                        \tag{11.1}
\end{align}
Here $\Delta$ is the source positive integer denoted
\texttt{Del}. At a top label its first Harish--Chandra
coordinate is $a_T(t)+(k-1)/2$. At a non-top label its first
coordinate $u_0$ and the adjacent coordinate $v_0$ satisfy
\[
     u_0-v_0\le C_cS_c,\qquad 0\le v_0\le S_c.
                                                        \tag{11.2}
\]
Every word member has the same first coordinate. These are the
equalities \texttt{top\_u}, \texttt{cloud\_gap},
\texttt{cloud\_v}, and \texttt{hcW\_first}; the highest
weight is obtained by subtracting the fixed $\rho$-vector.

\begin{proof}[Proof of Theorem~\ref{thm:highest}]
For $H\subset G$ of index at most eight, put $H_0=H\cap\Psi$.
Then $|H_0|\ge n/4$. Multiplication makes $H_0$ act freely
on $\Psi$, so every orbit has at least $n/4$ elements.
For sufficiently large $x$, $\delta<1/8$; (11.1) then says
$0<|\mathcal T|<n/4$. Thus $\mathcal T$ is not a union of
$H_0$-orbits. Some $h\in H_0$ moves a top label to a non-top label.

At a top label the first highest-weight coordinate is exactly
$a_T(t)\ge x\exp(x^\epsilon)$. At a non-top label it is
$u_0-(k-1)/2\le(C_c+1)S_c$. Put
$C_+=\max(1,C_c+1)$. Since $k,n\ge1$, the last inequality
in (11.1), which is a clause of \texttt{SizeDel}, gives
\[
 \frac{S_c}{a}\le\frac1{kxn\exp(x^{\epsilon/3})}<\frac1x.
                                                        \tag{11.3}
\]
For $x>C_+$, the first highest-weight coordinate at every top
label is strictly greater than at every non-top label. For all
sufficiently large $x$ the earlier cardinality condition also holds.
In particular the top set can be recovered from the
weight tuple itself, and that tuple is not $H_0$-invariant.

The source coefficient module, after complexification, is
\[
       W=\bigotimes_{\tau\in\Psi}V_{\lambda_\tau}
       \quad\text{for}\quad
       \mathcal G=\prod_{\tau\in\Psi}\mathrm{GL}_k(\mathbb C),
                                                        \tag{11.4}
\]
with the determinant shifts included. The irreducible modules
$V_{\lambda_\tau}$ have their indicated highest weights;
see \citet[Sections~20.3--20.8]{Milne2018}. The product
representation has a unique highest-weight line for the product
Borel. An intertwiner between it and the representation obtained
by permuting the factors by $h$ must carry this line to the
highest-weight line. Equality of their diagonal-torus characters
then forces $\lambda_{h\tau}=\lambda_\tau$ for every $\tau$.
The element found above violates this equality. Therefore no
action of all of $H_0$ can implement this coordinate permutation
on the one fixed coefficient representation.
\end{proof}

\subsection{What an equivariant enlargement would still require}

The conclusion concerns the natural field-permutation action,
not all possible linear actions. A direct sum over conjugated
weight tuples can formally carry an action permuting the
summands. This is an action relating different representations;
it does not by itself identify a commuting Galois action with
the Artin representation in (10.1), or a projected theta gain.
Also the reference complex embedding in the literal data can
change between transported objects. Although the family fixes
the labelled top set, this relabelling prevents an inference
that every conjugated object is absent from the family. The
proved obstruction is internal to one fixed labelled module.

There are already self-adjoint idempotents in the theta data.
Their specified finite group is a product of
$\mathrm{GL}_r(\mathbb F_{g_i})$ over the two auxiliary level
primes and their places, and the central idempotents select
cuspidal finite-level types. This group and action are distinct
from the proposed deck group $\Gal(F/F_0)$. The inspected
definitions supply no identification of them.

The column map is assembled from a theta unfolding map,
determinant characters, and products indexed by the embeddings,
with exponents fixed by the Harish--Chandra parameters. A deck
construction would have to transport these ingredients, the
Hermitian forms, class-lattice weights, integral lattices, and
descended theta and period maps coherently. Uniform valuation
inclusions would then pass as in (10.8), provided the projected
source is the one on which they hold. The exact prime-power
trace identity and its normalization would still have to be
identified. They do not follow from positivity alone.

Finally the literal centering term is also specific: its
\texttt{chiK} is the quadratic Dirichlet character modulo $d$,
and its period term uses
\[
 \log\Omega=\gamma+\log(\pi/v)
                     -\frac{L'(1,\chi_{-d})}{L(1,\chi_{-d})}.
                                                        \tag{11.5}
\]
Replacing this by a statistic of $J_0$ requires a new equality
of the local gain, archimedean terms, and centering constant.
The defining CM quadratic field of this literal construction
is $\Q(\sqrt{-d})$, not $K$. A construction over $K$ is a
modified theorem with its own hypotheses. The declaration that
all roots of unity are $\pm1$ is only an optional proposition in
the inspected host definitions; it is not used as a premise or
an obstruction in this paper.


\section{Scope and remaining alternatives}
The unconditional literal theorem excludes all ordinary small-norm
Hecke statistics, and arbitrary signed mixtures of them, on the odd
part of the comparison range. The common multiquadratic construction
shows that exact transfer and all five isolated numerical conditions
can coexist after the literal arithmetic conditions are relaxed.
The GRH discriminant and mixture estimates quantify restrictions in
that broader class.

The fixed-quotient projector is a complete representation-theoretic
construction with an exact integral cost. Its four-factor product is
a sufficient fixed analytic target. It has not been realized in the
theta chain. The highest-weight theorem excludes its simplest
coordinate-permutation implementation on one module. The remaining
possibility is a coherent enlarged construction with the correct
Frobenius traces, integral maps, archimedean normalization, and a
uniform gain on the relevant lattices. Such a uniform inclusion
inherits its valuation under a compatible projection; determinant
information alone does not. Nonabelian hosts under modified hypotheses and
comparisons involving the rational-prime orbit gain are also
outside the ordinary-trace exclusion. None follows from numerical
admissibility alone.

\section*{Note on method}
The text was produced by GPT-6 Astra in ChatGPT Pro mode. Separate
instances of the same model reconstructed the common-host construction
and the new literal-host and weighted-transfer statements from their
statements and the cited source definitions, before reading the
author's proofs. Their checks covered the empty norm range, exact
Frobenius orders, candidate and penalty conditions, the fixed-quotient
projector, its integral cost and ramified traces, the uniform mixture
bound, and the highest-weight obstruction. A further instance read
the complete manuscript. The uniform Dirichlet-character proof received a separate
targeted check by the same model. These are
distinct same-model checks; no human has checked the work.

Exact finite computations supplement the general proofs. They verify
three common-host constructions, forty-four one-prime fields,
4,365 odd prime-power exclusions, 1,824 projected trace identities,
and the stated projector indices and finite countermodels. They do
not verify GRH or a theta-chain realization. The formal source
definitions were inspected at the pinned revision; no Lean
compilation or verification of the surrounding development is claimed.
Precise check scopes, source locators, and hashes are recorded in the
public records. The text is dedicated to the public domain under
CC0~1.0. Editable material is at
\expandafter\url\expandafter{\workbenchurl}; the records are at
\expandafter\url\expandafter{\recordurl}.

\begin{thebibliography}{9}
\bibitem[Alp\"oge(2026a)]{Alpoge2026}
Alp\"oge, L. (2026a, November 20). \emph{Remark}.
\url{https://alpo.ge/remark.pdf}
\bibitem[Alp\"oge(2026b)]{AlpogeLean2026}
Alp\"oge, L. (2026b). \emph{Lean formalization accompanying Remark}
[Computer software]. Declarations \texttt{Frame}, \texttt{HostRaw},
\texttt{Host.rdE}, \texttt{Scale}, and \texttt{SizesAt}.
\url{https://github.com/CocoSgt/riemann-hypothesis-lean/tree/7cbf7a0c32fca1badf05931879a37c3f8a2b03b4}
\bibitem[Bach and Sorenson(1996)]{BachSorenson1996}
Bach, E., \& Sorenson, J. (1996). Explicit bounds for primes in
residue classes. \emph{Mathematics of Computation, 65}(216),
1717--1735. \url{https://doi.org/10.1090/S0025-5718-96-00763-6}
\bibitem[Milne(2018)]{Milne2018}
Milne, J. S. (2018). \emph{Reductive groups} (Version 2.00).
\url{https://www.jmilne.org/math/CourseNotes/RG.pdf}
\bibitem[Milne(2020)]{Milne2020}
Milne, J. S. (2020). \emph{Class field theory} (Version 4.03).
\url{https://www.jmilne.org/math/CourseNotes/CFT.pdf}
\bibitem[Poonen(2015)]{Poonen2015}
Poonen, B. (2015). \emph{Tate's thesis}. MIT 18.786 lecture notes.
\url{https://math.mit.edu/~poonen/786/notes.pdf}
\end{thebibliography}
\end{document}
